% ======================================================================
\section{The \texorpdfstring{\Htwo}{H2} level model}
\label{sec:h2model}
% ======================================================================

The level model computes, in every cell, the fractional populations of the
302 rovibrational levels of the electronic ground state $X\,^1\Sigma_g^+$ of
\Htwo{} under the cell's own radiation field, gas temperature and colliders,
and from them the Solomon photodissociation rate per molecule that the
chemistry uses (Section~\ref{sec:photorates}) and three \Htwo{} heating terms
(Section~\ref{sec:thermal}). The physics is in \src{h2\_mod.f90} (1570 lines);
the driver that supplies the field, the columns, the colliders and the
formation distribution is \rtn{h2LevelModel} with its internal routine
\rtn{solveCell} (\src{pdr\_mod.f90:\allowbreak{}5201--6014}). It is a re-implementation,
from the data and conventions of Cloudy's large \Htwo{} model
\citep{Shaw2005}, of the line-by-line treatment of ultraviolet pumping and
self-shielding that the Meudon PDR code \citep{LePetit2006} and the models of
\citet{DraineBertoldi1996} also use, in place of a shielding function fitted
to such models. It was validated level by level against Cloudy on one-zone and
slab models (\dlog{H-003}--\dlog{H-006}). Those comparisons isolate three
deliberate departures from Cloudy, each discussed below: the closure of the
level balance (Section~\ref{sec:balance}), the damping width of the lines
(Section~\ref{sec:shield}) and the treatment of line overlap
(Section~\ref{sec:overlap}).

\paragraph{Physical picture.} A molecule in level $l$ of $X$ absorbs a
far-ultraviolet photon in a Lyman ($B\,^1\Sigma_u^+$--$X$) or Werner
($C\,^1\Pi_u$--$X$) line, or in one of the weaker $B'\,^1\Sigma_u^+$--$X$ and
$D\,^1\Pi_u$--$X$ lines; out of $v=0$ the Lyman and Werner lines lie between 912
and 1110\,\AA{} (11.2--13.6\,eV), and the $B'$ and $D$ systems reach below
912\,\AA. Within about a nanosecond the excited level $u$ decays, either into
the vibrational continuum of $X$, which dissociates the molecule into two H(1s)
atoms (the Solomon process), or into a bound level of $X$, mostly
vibrationally excited, from which a slow electric-quadrupole cascade
($A=3\times10^{-11}$--$6\times10^{-6}$\,s$^{-1}$, median $2\times10^{-8}$\,s$^{-1}$ in
\src{transprob\_X.dat}) competes with collisions. In an unshielded field of the
Draine shape the model dissociates one molecule per 8.1 absorptions, at
$6.36\times10^{-11}$\,s$^{-1}$ per unit Draine field (\dlog{H-007}); the mean
dissociation probability per absorption is 0.12 in the computations of
\citet{Sternberg2014} and about 0.15 when averaged through complete
photodissociation fronts \citep{DraineBertoldi1996}. The strong lines out of a
level become optically thick at level columns of order
$10^{14}$\,cm$^{-2}$, saturate, and above $\sim10^{17}$--$10^{18}$\,cm$^{-2}$ absorb
mainly in their Lorentzian wings, where above $N(\mathrm{H_2})\sim10^{20}$\,cm$^{-2}$
the lines of different levels overlap \citep{DraineBertoldi1996}. The
dissociation rate in a cell therefore depends on the field at every line, on
the column of every level along the paths of the photons, on the line widths,
and on the populations of the absorbing levels, which the pumping itself
changes. The model treats all of these line by line.

\subsection{Data}
\label{sec:h2data}

All data are Cloudy c25.00's (\src{dataPDR/\allowbreak{}h2/\allowbreak{}}, copied
unmodified; provenance in \src{dataPDR/\allowbreak{}h2/\allowbreak{}README.md}). The
reading rules are Cloudy's: a line starting with \texttt{\#} is a comment, and the
data of the excited-state files end at the first empty line. The sources in
Table~\ref{tab:h2data} are those that the files' own headers give, completed
where the header is not the primary source.

\begin{longtable}{@{}L{0.33\linewidth}L{0.28\linewidth}L{0.33\linewidth}@{}}
\caption{Data of the \Htwo{} level model. Temperature ranges are those of the
tables; coverage is given as the upper levels present.}\label{tab:h2data}\\
\toprule
Content & File(s) & Source \\
\midrule
\endfirsthead
\toprule
Content & File(s) & Source \\
\midrule
\endhead
energies of the $X$ rovibrational levels (302 levels, $v\le14$, $J\le31$; sorted by energy) & \src{energy\_X.dat} & \citet[their Table~1]{Komasa2011}, compiled by K.~Kaplan (header) \\
radiative transition probabilities within $X$: electric quadrupole, plus magnetic-dipole components of 30 $v=1\to0$, $\Delta J=0$ transitions listed separately and summed & \src{transprob\_X.dat} & \citet{Wolniewicz1998}; magnetic dipole \citet{PachuckiKomasa2011} \\
H\,+\,\Htwo{}: 1432 transitions, $v\le3$, $J\le17$, 728 of them ortho$\leftrightarrow$para (reactive); 100--5000\,K & \src{coll\_rates\_H\_15.dat} & \citet{Lique2015} \\
He\,+\,\Htwo{}: 3251 transitions, $v\le14$, $J\le30$, no spin change; 1--$10^4$\,K & \src{coll\_rates\_He\_ORNL.dat} & ORNL data set; the header cites ``Lee, T.\ G.\ et al.\ 2007, ApJ, in preparation''; no published source for this table was identified \\
ortho- and para-\Htwo\,+\,\Htwo{}: 240 and 241 transitions, $v=0$, $J\le31$; 1--$10^4$\,K & \src{coll\_rates\_H2ortho\_ORNL.dat}, \src{coll\_rates\_H2para\_ORNL.dat} & \citet{Wan2018} \\
\Hp\,+\,\Htwo{}: 44 transitions, $v=0$, $J\le9$, 24 of them ortho$\leftrightarrow$para (proton exchange); 3--6000\,K & \src{coll\_rates\_Hp.dat} & \citet{Gerlich1990} \\
energies of the rovibrational levels of $B$, $C^+$, $C^-$, $B'$, $D^+$, $D^-$ (1592 levels) & \src{energy\_\{B,C\_plus,\ldots\}.dat} & derived from \citet{Abgrall2000} (headers) \\
transition probabilities of the 32\,453 lines to $X$ & \src{transprob\_\{\ldots\}.dat} & \citet{Abgrall2000}, from the band-system calculations of \citet{Abgrall1993} ($B$--$X$, $C$--$X$) and \citet{Abgrall1994} ($B'$--$X$, $D$--$X$) \\
spontaneous radiative dissociation rate $D_u$ [s$^{-1}$] and mean kinetic energy of the atoms [eV] per excited level & \src{dissprob\_\{\ldots\}.dat} & \citet{Abgrall2000} \\
dissociation energy of $X$, $D=36\,118.11$\,cm$^{-1}$ (4.478\,eV; the only value of the file the model uses) & \src{energy\_dissoc.dat} & \citet{Sharp1970} per the header (which dates it 1971), with values revised in Cloudy in 2002 \\
\bottomrule
\end{longtable}

Reading: \rtn{h2ReadX} (\src{h2\_mod.f90:\allowbreak{}206--334}) and \rtn{h2ReadElec}
(\src{:\allowbreak{}366--516}). Statistical weights are $g=(2J+1)(2I+1)$, with
nuclear-spin factor $2I+1=3$ for ortho and 1 for para. In $X$ ortho levels have
odd $J$; in the excited states $B$, $B'$ and the $\Pi^+$ components $C^+$, $D^+$
ortho levels have even $J$, and in the $\Pi^-$ components $C^-$, $D^-$ odd $J$,
following Cloudy \citep{Shaw2005}: this is the assignment under which every
allowed transition ($\Delta J=\pm1$ to the $\Sigma_u^+$ states and the $\Pi^+$
components, $\Delta J=0$ to the $\Pi^-$ components) conserves the nuclear spin,
and the reader counts the lines that would not (none, \dlog{H-004}). The $\Pi$
states have no $J=0$ level. The $D^-$ state stops at $v=2$ because Cloudy's file
has an empty line before the levels it marks ``above the dissociation limit'';
3138 $D^-$ lines are dropped with them, and 392 lines whose lower level lies
beyond the 302 levels of $X$ (\dlog{H-004}). The dissociation rates $D_u$ are
those of spontaneous emission into the vibrational continuum of $X$
\citep{Abgrall2000}. Apart from 45 high-$J$ levels of $D^-$ ($v\le2$), all the
excited levels lie below the H(1s)\,+\,H($n=2$) limit (118\,376\,cm$^{-1}$ in
\src{energy\_dissoc.dat}), where this radiative channel is the only way they
dissociate; predissociation, which acts above that limit, is not in the data.
A transition listed
twice adds its components, as Cloudy does; reading the $X$ file without adding
the magnetic-dipole components had doubled the $v=1$ populations at low density
(\dlog{H-003}). The lines are stored sorted by upper level (counting sort) and
indexed by wavenumber (heapsort) for the overlap calculation.

\subsection{Statistical equilibrium and the formation closure}
\label{sec:balance}

Let $n_i$ be the fraction of the cell's \Htwo{} molecules in level $i$ of $X$
($\sum_i n_i=1$). The populations solve (\rtn{balance},
\src{h2\_mod.f90:\allowbreak{}1456--1477})
\begin{equation}
  \sum_{j\neq i} n_j R_{ji} - n_i\Bigl(\sum_{j\neq i} R_{ij} + k_i\Bigr)
  + \phi_i \sum_j n_j k_j + s_i = 0,\qquad i=1,\ldots,302,
  \label{eq:balance}
\end{equation}
where $R_{ij}$ [s$^{-1}$] is the total rate per molecule from level $i$ to
level $j$ (collisions, spontaneous decay within $X$ and the fluorescent
transfer through the excited states; Sections~\ref{sec:coll} and
\ref{sec:pump}), $k_i$ [s$^{-1}$] the Solomon dissociation rate out of level
$i$, $\phi_i$ the fraction of newly formed molecules that enter level $i$
($\sum_i\phi_i=1$; Section~\ref{sec:formdist}) and $s_i$ [s$^{-1}$] an optional
further source per molecule, not used by \rtn{solveCell}. Fluorescent returns
to the absorbing level ($R_{ii}$) drop out. Summed over $i$, the transfer terms
cancel and so do the dissociation and formation terms, so the system
conserves $\sum_in_i$: the term $\phi_i\sum_jn_jk_j$ is the closure \emph{every
molecule destroyed by the Solomon process is replaced by one formed}. The
level model determines how the molecules are distributed over the levels;
how many there are is the chemistry's (Section~\ref{sec:chemistry}).

\paragraph{Numerics.} The equation of level 1 is replaced by $\sum_in_i=1$
and the $302\times302$ system is solved directly by Gaussian elimination with
partial pivoting (\rtn{gaussSolve}, \src{:\allowbreak{}1536--1568}; about
$n^3/3\approx9\times10^6$ operations per solution). \rtn{h2Solve}
(\src{:\allowbreak{}1516--1533}) assembles $R=C+A+W$ (collisions, radiative decay,
fluorescent transfer) and the $k_i$, and calls it. The populations of the
validation runs were never negative and summed to unity to the printed
precision (\dlog{H-003}). Cloudy instead solves its lowest 70 levels by matrix
inversion and the others one by one (\dlog{H-010}).

\paragraph{Validity of the closure.} In steady state the chemistry's formation
rate per molecule, $R_{\mathrm f}\,n_{\mathrm H}\,n(\mathrm H)/n(\mathrm{H_2})$, equals
the sum of all destruction rates per molecule, of which $k_{\mathrm{Sol}}$ is
one; the closure sets the formation flux into the levels equal to
$k_{\mathrm{Sol}}$ alone. It is therefore exact where photodissociation
dominates the destruction (93--100\% of the \Htwo{} destruction out to 0.43\,pc
in \texttt{pn150k}, \dlog{H-009}), and underestimates the formation flux by
the factor $k_{\mathrm{Sol}}/(k_{\mathrm{Sol}}+k_{\mathrm{other}})$ elsewhere. The
network's three \Htwo\,+\,CRP channels alone give $k_{\mathrm{other}}=0.99\,\zeta$,
$1.0\times10^{-16}$\,s$^{-1}$ at the default \kw{pdrZeta}, so the closure fails
where the Solomon rate is shielded below a few $10^{-16}$\,s$^{-1}$, a reduction by
more than $10^5$ from its unshielded value at $\chi\sim1$. There the formation
pumping of the $v\ge1$ levels, the infrared emission it produces and the part
of the formation energy that reaches the gas by collisional de-excitation are
underestimated; the Solomon rate itself is affected only through the
populations of the absorbing levels, and the kinetic part of the formation
energy follows the chemistry's own formation rate
(Section~\ref{sec:thermal}; Section~\ref{sec:limitations}). Cloudy closes its
level system differently: it applies its network's Solomon rate a second time
as a uniform sink on the levels below 4100\,cm$^{-1}$, solves without a sum
rule and rescales all populations by one factor (1.06--1.11 in one-zone models
with $n_{\mathrm H}\le10^4$\,cm$^{-3}$ and fields up to 100 Habing). Put through that
scheme, this model's rates reproduce Cloudy's populations to 0.002--0.018\,dex
in 15 of 16 such models, against a median difference of 0.07\,dex with its own
closure (\dlog{H-004}); Cloudy's closure raises the pumped line intensities by
about 1.2 (\dlog{H-009}).

\subsection{Collisional and radiative rates within \texorpdfstring{$X$}{X}}
\label{sec:coll}

\rtn{h2CollRad} (\src{h2\_mod.f90:\allowbreak{}1369--1415}) adds, for every pair of
levels $u>l$ (energies $E$ in cm$^{-1}$),
\begin{equation}
  R_{ul} \mathrel{+}= \sum_c n_c\,k^c_{ul}(T) + A_{ul},\qquad
  R_{lu} \mathrel{+}= \sum_c n_c\,k^c_{ul}(T)\,\frac{g_u}{g_l}\,e^{-hc(E_u-E_l)/kT},
  \label{eq:collrad}
\end{equation}
with $n_c$ [cm$^{-3}$] the density of collider
$c\in\{\mathrm H,\,\mathrm{He},\,o\text{-H}_2,\,p\text{-H}_2,\,\mathrm H^+\}$, $k^c_{ul}$
[cm$^3$\,s$^{-1}$] its de-excitation rate coefficient, $A_{ul}$ [s$^{-1}$] the
radiative transition probability, and $T=\max(T_e,1\,\mathrm K)$
(\src{pdr\_mod.f90:\allowbreak{}5667}):
\begin{itemize}
\item downward coefficients from the tables of Table~\ref{tab:h2data},
  interpolated linearly in $T$ and held at the end values outside the
  tabulated range (\rtn{tabRate}, \src{h2\_mod.f90:\allowbreak{}642--660}; Cloudy's
  \texttt{InterpCollRate} convention). Below 100\,K the H\,+\,\Htwo{} rates are
  therefore those of 100\,K, which overestimates the reactive (ortho--para)
  channels, whose barrier makes them fall steeply with temperature; above
  5000\,K (H), 6000\,K (\Hp) and $10^4$\,K (He, \Htwo) the rates are those of the
  table's end;
\item for a pair of levels of the same nuclear spin that a collider's table
  does not contain, the ``g-bar'' estimate of \citet[their eq.~1 and
  Table~2]{Shaw2005},
  \begin{equation}
    \log_{10}\bigl(k^c_{ul}/\mathrm{cm^3\,s^{-1}}\bigr) = y_0 + a_g\,\bigl[\max(\Delta E,\,100)\bigr]^{b_g},
    \qquad \Delta E = E_u-E_l\ [\mathrm{cm^{-1}}],
    \label{eq:gbar}
  \end{equation}
  a fit, per collider, of the rate coefficients then available as a function
  of the energy gap (coefficients in Table~\ref{tab:gbar}; Cloudy's
  \texttt{GbarRateCoeff}), which \citet{Shaw2005} found to lie within 1\,dex of
  60--70\% of the rates it was fitted to. A pair that changes the nuclear spin
  and is not in a table has no collisional rate;
\item upward rates by detailed balance at the gas temperature, so that
  collisions alone relax the populations to a Boltzmann distribution at $T$;
\item spontaneous radiative decay $A_{ul}$ (quadrupole plus magnetic dipole)
  from \src{transprob\_X.dat}.
\end{itemize}

\begin{table}[H]
\centering
\caption{g-bar coefficients of Equation~(\ref{eq:gbar}) (\src{h2\_mod.f90:\allowbreak{}96--101}),
identical to Table~2 of \citet{Shaw2005}.}\label{tab:gbar}
\begin{tabular}{@{}lrrr@{}}
\toprule
Collider & $y_0$ & $a_g$ & $b_g$ \\
\midrule
H & $-9.9265$ & $-0.1048$ & 0.456 \\
He & $-8.281$ & $-0.1303$ & 0.4931 \\
$o$-\Htwo & $-10.0357$ & $-0.0243$ & 0.67 \\
$p$-\Htwo & $-8.6213$ & $-0.1004$ & 0.5291 \\
\Hp & $-9.2719$ & $-0.0001$ & 1.0391 \\
\bottomrule
\end{tabular}
\end{table}

\paragraph{Coverage.} Of the 22\,675 pairs of levels of the same nuclear spin,
the tables contain 703 for H, 3250 for He, 240 and 239 for $o$- and $p$-\Htwo{}
and 20 for \Hp{} (counted from the data files); the remaining 86--99.9\% take
Equation~(\ref{eq:gbar}). Ortho and para are coupled by collisions only through
the 728 reactive H\,+\,\Htwo{} transitions ($v\le3$) and the 24 proton-exchange
\Hp\,+\,\Htwo{} transitions ($v=0$, $J\le9$), the latter being by far the faster
\citep{Shaw2005}. The tabulated pairs include the low-$J$ levels of $v=0$ that
carry most of the column; the g-bar rates matter where collisions compete
with radiative decay in levels without data, i.e.\ high-$v$ or high-$J$ levels
at high density, and \citet{Shaw2005} found the populations of the tabulated
low levels nearly insensitive to them.

\paragraph{Colliders.} The collider densities (\src{pdr\_mod.f90:\allowbreak{}5887--5888})
are taken from the chemistry's abundances of the previous iteration,
$n(\mathrm H)=n_{\mathrm H}x(\mathrm H)$, $n(\mathrm{He})$ and $n(\mathrm H^+)$, and
from the cell's own solution for the spin split,
$n(o\text{-H}_2)=f_o\,n(\mathrm{H_2})$, $n(p\text{-H}_2)=(1-f_o)\,n(\mathrm{H_2})$, with
$f_o=\sum_{i\in\mathrm{ortho}} n_i$ iterated to self-consistency
(Section~\ref{sec:solvecell}). Electrons are not a collider, as in Cloudy;
\citet{AlemanGruenwald2011} include them in their models of the \Htwo{}
emission of the ionised region of planetary nebulae. Where electrons are a
large fraction of the colliders, in the ionised gas and at the ionisation
front, the collisional excitation of \Htwo{} is therefore underestimated, in
gas whose \Htwo{} abundance is in any case uncertain (\dlog{L-43}). The rate
coefficients and Boltzmann factors are cached for the last temperature
(\rtn{crKd}, \rtn{crBf}), so the iterations over the colliders' densities at
fixed $T$ rebuild only $R$.

\subsection{Ultraviolet pumping, Solomon dissociation and trapping}
\label{sec:pump}

A molecule in $X$ level $l$ is excited in the line $l\to u$ at the rate
\begin{equation}
  P_{lu} = B_{lu}\,\bar J_{lu} = A_{ul}\,\frac{g_u}{g_l}\,\eta_{lu}\quad[\mathrm{s^{-1}}],
  \qquad
  \eta_{lu} = \frac{c^2}{2h\nu_{lu}^3}\,J_\nu(\nu_{lu}),
  \label{eq:pump}
\end{equation}
where $\bar J_{lu}$ is the mean intensity averaged over the line profile and
$\eta_{lu}$ the photon occupation number at the line after shielding
(Sections~\ref{sec:eta} and \ref{sec:shield}); the second form follows from
the Einstein relations for a field constant across the line. In the
unshielded Draine field $\eta\approx2\times10^{-19}$ near 12\,eV.

The excited levels have radiative lifetimes of 0.5--9\,ns in these data
(total decay rates $\Gamma_u=1.1\times10^8$--$1.9\times10^9$\,s$^{-1}$, median
$\approx10^9$\,s$^{-1}$ for $B$ and $C$). Collisions cannot compete with that below
$\sim10^{14}$\,cm$^{-3}$ \citep{Shaw2005}, and the excited populations, of order
$\sum_lP_{lu}/\Gamma_u\sim10^{-18}\chi$ per molecule, are never a reservoir: every
absorption is followed at once by a decay, into the continuum of $X$
(dissociation) or into a bound level of $X$. The excited states are therefore
not solved for. With the effective decay rates
\begin{equation}
  A^{\mathrm{eff}}_t = A_t\,(1+\eta_t)\,\beta_t,
  \qquad
  \Gamma^{\mathrm{eff}}_u = D_u + \sum_{t\in u} A^{\mathrm{eff}}_t ,
  \label{eq:aeff}
\end{equation}
where $t$ runs over the lines from $u$ to bound levels of $X$, $1+\eta_t$
accounts for stimulated emission (negligible, kept for completeness), $\beta_t$
is the escape probability of a photon emitted in line $t$
(Section~\ref{sec:escape}) and $D_u$ the spontaneous radiative dissociation rate
of $u$, \rtn{h2PumpRates} (\src{h2\_mod.f90:\allowbreak{}540--572}) returns
\begin{align}
  k_l &= \sum_u P_{lu}\,\frac{D_u}{\Gamma^{\mathrm{eff}}_u}
       && \text{(Solomon dissociation out of } l\text{)},\\
  W_{ll'} &= \sum_u P_{lu}\,\frac{A^{\mathrm{eff}}_{ul'}}{\Gamma^{\mathrm{eff}}_u}
       && \text{(fluorescent transfer } l\to l'\text{)},\\
  H^{\mathrm{dis}}_l &= \sum_u P_{lu}\,\frac{D_u}{\Gamma^{\mathrm{eff}}_u}\,E^{\mathrm{kin}}_u
       && \text{(kinetic energy given to the atoms [eV\,s}^{-1}\text{])},
\end{align}
and the total absorption rate $p_l=\sum_u P_{lu}$, with $E^{\mathrm{kin}}_u$ the
mean kinetic energy released per dissociation of $u$ (\src{dissprob\_\ldots.dat}).
The first step of the fluorescent cascade is $W$; the quadrupole cascade that
follows is carried by $A$ in Equation~(\ref{eq:balance}). The Solomon rate per
molecule handed to the chemistry is
\begin{equation}
  k_{\mathrm{Sol}} = \sum_l n_l\,k_l
  \qquad(\text{\src{pdr\_mod.f90:\allowbreak{}5918}}).
  \label{eq:ksol}
\end{equation}

\paragraph{Trapping.} A photon emitted in a line $t=u\to l'$ that is optically
thick is absorbed again close by, by a molecule in $l'$, which it re-excites
to the same level $u$. In the escape-probability (``on-the-spot'')
approximation the net effect is that the decay rate of the line is reduced to
$A_t\beta_t$, and the level's other decays and its dissociation take its share
in $\Gamma^{\mathrm{eff}}_u$. Only lines that end in populated low levels are
thick, so trapping closes the decays to $v=0$ and raises the dissociation
probability $D_u/\Gamma^{\mathrm{eff}}_u$ in thick gas. It is needed: without it
the rate into $v=0$, $J=0$--5 from the excited states was up to 2200 times
Cloudy's in the shielded zones of the slab of \dlog{H-005} and the Solomon rate
10--20\% low; with it the Solomon rate comes within 3--8\% of Cloudy's
(\dlog{H-005}). The total absorption rate $p_l$ counts only the external field.
Photons that dust or an overlapping \Htwo{} line absorb before they escape are
not counted as lost (Cloudy adds such a destruction probability to the escape
probability, \citealt{Shaw2005}); this slightly overestimates the dissociation
probability of levels whose decays are trapped.

\subsection{Self-shielding of the pumping lines}
\label{sec:shield}

\paragraph{Line parameters.} For line $lu$ of wavelength $\lambda=1/\tilde\nu_{lu}$
[cm] (\rtn{h2LineParams}, \src{h2\_mod.f90:\allowbreak{}1328--1336}),
\begin{equation}
  \sigma_{lu} = \frac{g_u}{g_l}\,\frac{A_{ul}\,\lambda^3}{8\pi^{3/2}\,b}\ [\mathrm{cm^2}],\qquad
  \tau_0 = \sigma_{lu}N_l,\qquad
  a = \frac{\Gamma_u\,\lambda}{4\pi\,b},\qquad
  \Gamma_u = D_u + \sum_{t\in u}A_t ,
  \label{eq:lineparams}
\end{equation}
where $\sigma_{lu}$ is the line-centre cross section of a pure Doppler profile
of Doppler width $\Delta\nu_D=b/\lambda$, $(\pi e^2/m_ec)f_{lu}/(\sqrt\pi\Delta\nu_D)$
written with $A_{ul}$; $N_l$ [cm$^{-2}$] is the column of molecules in level $l$;
$b=\sqrt{2kT/m_{\mathrm{H_2}}+v_{\mathrm{turb}}^2}$ [cm\,s$^{-1}$] the Doppler parameter of
the target cell (\src{pdr\_mod.f90:\allowbreak{}5668}; $v_{\mathrm{turb}}$ is \kw{pdrVturb},
default 1\,km\,s$^{-1}$); and $a$ the Voigt damping parameter, the ratio of the
Lorentzian half width $\Gamma_u/4\pi$ to $\Delta\nu_D$. $\Gamma_u$ is the total
decay rate of the upper level, bound--bound plus bound--free (spontaneous
radiative dissociation), i.e.\ the inverse radiative lifetime; since the lower
level of every line is a metastable level of $X$, $\Gamma_u$ is the natural width
of the line. With these data $a\le0.03$ for $b\ge1$\,km\,s$^{-1}$ (median 0.006 at
100\,K and the default $v_{\mathrm{turb}}$).

\paragraph{The damping width.} $\Gamma_u$ exceeds a line's own $A_{ul}$ by a
factor 20 in the median over the thick lines of the \dlog{H-005} slab (16--84\%:
7.8--157), and the shielding in the damping wings scales as $a^{1/2}$
(Equation~\ref{eq:Sasym}). Cloudy damps each \Htwo{} line with its own $A_{ul}$
(\texttt{dampXvel}$\,=A_{ul}\lambda/4\pi$ in \texttt{mole\_h2\_create.cpp} of c25.00,
and in its saved damping constants, \dlog{H-005}), whereas its hydrogenic lines
use the lifetime of the upper level (\texttt{iso\_create.cpp}). The natural width
is the physical choice. Cloudy's lowers the transmission in the damping regime
by up to $\sqrt{20}\approx4.5$ (pumped lines behind more than $10^{16}$\,cm$^{-2}$ of
\Htwo{} fall to 0.18--0.41 of their intensity, \dlog{H-008}) and moves the
H/\Htwo{} crossing 0.21\,mag shallower on \texttt{hiipdr} and 0.10\,mag on
\texttt{pn150k} (\dlog{H-026}; Section~\ref{sec:limitations}). The same $b$ (the target cell's) is used for the
whole column in front of the cell and along the escape rays; this matters in
the Doppler-saturated regime, where $\tau_0\propto b^{-1}$, and not in the damping
regime, where the equivalent width in frequency does not depend on $b$.

\paragraph{Transmission and curve of growth.} The rate at which a molecule
behind the column $N_l$ absorbs in the line, relative to the unshielded rate
for a continuum flat across the line, and its integral over $\tau_0$ are
\begin{align}
  S(\tau_0,a) &= \frac{1}{\sqrt\pi}\int_{-\infty}^{\infty} H(a,x)\,
       \exp\!\left[-\tau_0\frac{H(a,x)}{H(a,0)}\right]\dd x,
  \label{eq:SW}\\
  W(\tau_0,a) &= \int_0^{\tau_0} S\,\dd\tau
      = \frac{2H(a,0)}{\sqrt\pi}\int_0^{\infty}\left\{1-\exp\!\left[-\tau_0\frac{H(a,x)}{H(a,0)}\right]\right\}\dd x ,
  \label{eq:W}
\end{align}
with $x=(\nu-\nu_{lu})/\Delta\nu_D$ and $H(a,x)=(a/\pi)\int e^{-y^2}[(x-y)^2+a^2]^{-1}\dd y$
the Voigt function ($\int H\,\dd x=\sqrt\pi$), evaluated with the rational
approximation of \citet{Humlicek1982} (relative accuracy $10^{-4}$; \rtn{h2Voigt},
\src{h2\_mod.f90:\allowbreak{}664--686}). $S$ is the single-line self-shielding
function: the profile-weighted transmission, $S(0,a)=1$. $W$ is the equivalent
width in units of $\sqrt\pi\Delta\nu_D/H(a,0)$, normalised so that $W\to\tau_0$ for
$\tau_0\ll1$, and $\dd W/\dd\tau_0=S$. The three regimes are
\begin{equation}
  S \simeq \begin{cases}
    1-\tau_0/\sqrt2 & (\tau_0\ll1),\\[2pt]
    \bigl[\sqrt\pi\,\tau_0\,(\ln\tau_0)^{1/2}\bigr]^{-1} & \text{(saturated Doppler core)},\\[2pt]
    \pi^{-1/4}\,\bigl[H(a,0)\,a/\tau_0\bigr]^{1/2} & \text{(damping wings)},
  \end{cases}
  \label{eq:Sasym}
\end{equation}
the first two for $a\to0$, the third when the absorption is dominated by the
Lorentzian wings $H\simeq a/(\sqrt\pi x^2)$.

\paragraph{Profile normalisation.} $\tau_0=\sigma_{lu}N_l$ is the line-centre
optical depth of the Doppler core, and $S$ normalises the profile to its own
centre, $\tau(x)=\tau_0H(a,x)/H(a,0)$; with this convention $\tau_0$ agrees with
Cloudy's line optical depths to 0.04\% (\dlog{H-005}). The area-normalised Voigt
profile gives $\tau(x)=\tau_0H(a,x)$, smaller by the factor
$H(a,0)\simeq1-2a/\sqrt\pi$ at every frequency. For these data the difference is
at most 3\% in $\tau$ (0.8\% weighted by the unshielded pumping, and half that in
$S$ in the damping regime), negligible against the other uncertainties.

\paragraph{Numerical strategy.} \rtn{shieldQuad} (\src{:\allowbreak{}692--746})
integrates one side of the profile with the trapezoidal rule on 1200 uniform
points to $x=6$ and logarithmically (60 points per decade) beyond, to
$x_{\max}=\max(60,\,300\sqrt{\tau_0 a})$; beyond $x_{\max}$ the Lorentzian wing
$H\simeq a/(\sqrt\pi x^2)$ is optically thin ($\tau\le10^{-5}$ at $x_{\max}=300\sqrt{\tau_0a}$)
and is added analytically. It also records the offsets $x_q$ at which $1/6$,
$1/2$ and $5/6$ of one side's absorption, $\int H e^{-\tau}\dd x$, has accumulated,
for the overlap quadrature. Against an independent quadrature (the Faddeeva
function of SciPy) $S$ and $W$ agree to 0.04\% and 0.03\% over
$\tau_0=10^{-2}$--$10^{11}$, $a=10^{-6}$--$10^{-1}$ (\dlog{H-005}). \rtn{shieldTable}
(\src{:\allowbreak{}748--760}) tabulates $\log S$, $\log W$ and $\log x_q$ once on a
$161\times29$ grid, $\log_{10}\tau_0\in[-4,16]$ (step 0.125) and
$\log_{10}a\in[-7,0]$ (step 0.25). Look-up (\rtn{tabShieldJ}, \src{:\allowbreak{}833--851})
is bilinear in $(\log\tau_0,\log a)$ on the logarithms; below the table $S=1$ and
$W=\tau_0$; above it the damping-wing asymptote, $S\propto\tau_0^{-1/2}$,
$W\propto\tau_0^{1/2}$. The table's $S$ is within 1.2\% of the exact single-line
integral (\dlog{H-006}). The damping parameter's column and fraction in the table
are found once per line (\rtn{h2ShieldAIndex}, \src{:\allowbreak{}822--830}).

\paragraph{Comparison with shielding functions.} A shielding-function
treatment writes the Solomon rate as $k_{\mathrm{Sol}}\simeq k_{\mathrm{Sol}}(0)\,
f_{\mathrm{shield}}(N_{\mathrm{H_2}})\,e^{-\tau_{d,1000}}$ \citep[their
eq.~40]{DraineBertoldi1996}, with $\tau_{d,1000}$ the dust optical depth at
1000\,\AA{} and $f_{\mathrm{shield}}$ a fit to multi-line models for given
populations, temperature and field shape. The fit of \citet[their
eq.~37]{DraineBertoldi1996},
\begin{equation}
  \begin{split}
  f_{\mathrm{shield}} &= \frac{0.965}{(1+x/b_5)^2}+\frac{0.035}{(1+x)^{1/2}}\,
     \exp\!\left[-8.5\times10^{-4}(1+x)^{1/2}\right],\\
  x&=\frac{N(\mathrm{H_2})}{5\times10^{14}\,\mathrm{cm^{-2}}},\qquad b_5=\frac{b}{\mathrm{km\,s^{-1}}},
  \end{split}
  \label{eq:db96}
\end{equation}
describes the saturation of the lines in its first two terms and their
overlap in the exponential. 3D-PDR's own rate uses the single-line function of
\citet{Federman1979} for one representative line
(Section~\ref{sec:h2rate}). Here the equivalent of $f_{\mathrm{shield}}$
follows from the level columns, the line widths, the populations and the
field at each line, and the dust attenuation is in the MC field. On the Cloudy
slab of \dlog{H-005} (100\,K, $n_{\mathrm H}=10^4$\,cm$^{-3}$, 100 Habing,
$b=0.91$\,km\,s$^{-1}$) the model without overlap followed Equation~(\ref{eq:db96})
without its overlap factor to 0.79--1.00 at $N(\mathrm{H_2})=10^{18}$--$10^{21}$\,cm$^{-2}$,
and with overlap it lies at 0.70--0.95 of the full fit there and 20--50\% above
it at $10^{14}$--$10^{17}$\,cm$^{-2}$ (\dlog{H-005}, \dlog{H-006}). The analytic
theory of the H\,\textsc{i}-to-\Htwo{} transition of \citet{Sternberg2014}, in
which the total atomic column of a slab in a beamed field is
$N_{\mathrm{H\,I}}=\sigma_g^{-1}\ln(\alpha G/2+1)$ (their eq.~40), with $\sigma_g$ the
far-ultraviolet dust absorption cross section per H nucleus, $\alpha$ the ratio
of the free-space dissociation rate to the formation rate $R_{\mathrm f}n$ and $G$ a
cloud-averaged self-shielding factor, rests on the same kind of line-by-line
computation (an ``effective dissociation bandwidth'' computed with the Meudon
code).

\subsection{Averaging through the cell's own column}
\label{sec:ownavg}

The pumping rate that the level balance of a cell needs is its mean over the
cell, whose populations are uniform. If the cell adds $\Delta\tau=\sigma_{lu}\Delta N_l$
to the column $\tau_0$ in front of it, the mean of $S$ over a uniform layer
traversed along one chord is the mean slope of the curve of growth across it
(\rtn{h2Shield}, \rtn{h2ShieldJ}, \src{h2\_mod.f90:\allowbreak{}768--784, 854--866};
\rtn{h2LineShieldPrep/Eval}, \src{:\allowbreak{}876--907}):
\begin{equation}
  \langle S\rangle = \frac{W(\tau_0+\Delta\tau,a)-W(\tau_0,a)}{\Delta\tau}
  \quad\text{if }\Delta\tau > 10^{-3}\max(\tau_0,10^{-3}),
  \qquad\text{else } S(\tau_0+\tfrac12\Delta\tau,a),
  \label{eq:cellavg}
\end{equation}
clamped to $[0,1]$; the second form avoids the cancellation in the difference
of $W$ for thin layers, with an error of order $\Delta\tau^2 S''/24$. A cell whose own
column takes its lines from thin to saturated absorbs the band in its front
part, and $\langle S\rangle$ is then much smaller than $S(\tau_0)$. The cell's own
column of level $l$ is $\Delta N_l = n_l\,n(\mathrm{H_2})\,\ell$ with two path
lengths (\src{pdr\_mod.f90:\allowbreak{}5609--5615, 5697--5706, 5804--5819}):
\begin{align}
  \ell_{\mathrm{dir}} &= \frac{\Delta x\,\Delta y\,\Delta z}{\Delta y\Delta z\,|u_x|+\Delta x\Delta z\,|u_y|+\Delta x\Delta y\,|u_z|},\\
  \ell_{\mathrm{iso}} &= \frac{4V}{A} = \frac{2\,\Delta x\,\Delta y\,\Delta z}{\Delta x\Delta y+\Delta y\Delta z+\Delta x\Delta z},
\end{align}
where $\ell_{\mathrm{dir}}$ is the mean chord of the cell along the unit vector
$\hat u$ towards the source (the volume over the area projected on a plane
normal to $\hat u$), averaged over sources with weights $r_s^{-2}$ and used for
the direct beam, and $\ell_{\mathrm{iso}}$ is Cauchy's mean chord of a convex body
for isotropically distributed lines, used for the diffuse light. The widths
$\Delta x,\Delta y,\Delta z$ are centred differences of the axis values
(\rtn{cWid}), one-sided at the first and last axis value; there the full
spacing is used, although the outermost cell of a grid that ends at its last
axis value is half as wide. One representative chord stands for the
distribution of chord lengths through the cell. The direct-path column uses
the populations of the current pass of the cell's iteration, the diffuse one
those of the previous sweep.

This is how the level model includes the self-shielding of the cell's own
\Htwo, which the chemistry's columns leave out (Section~\ref{sec:avradial}).
Averaging the transmission through the cell is a different operation from
adding the cell's near half to the shielding column, which made the
chemistry's surface cell shield itself completely and settle on a
self-shielded branch (\dlog{D-088}). The size of the effect shows how thin the
pumping layer is: when the diffuse part was first computed without the cell's
own column, the Solomon rate where \Htwo{} turns on was 4--12 times the
chemistry's, and averaging through the mean chord lowered the 1--0\,S(1)
luminosity of the test model by a factor 4.6 (\dlog{H-007}). The
average is exact for a uniform cell; cells across the H/\Htwo{} transition are
not, and on \texttt{hiipdr} the 1--0\,S(1) luminosity changed by $+35\%$ from
$21^3$ to $31^3$ and by $+3\%$ from $31^3$ to $41^3$ cells while the \Htwo{} profile
of the transition was still moving (\dlog{H-007}).

\subsection{Line overlap}
\label{sec:overlap}

The \Htwo{} lines shade one another. Each line's transmission is its own $S$
times the mean transmission of all the \emph{other} lines over the frequencies
where it absorbs (\rtn{lineOverlap}, \src{h2\_mod.f90:\allowbreak{}1018--1289},
\dlog{H-006}):
\begin{equation}
  \mathcal{O}_i = \frac16\sum_{k=1}^{3}\sum_{s=\pm1}
     \exp\!\Bigl[-\tau_{\mathrm{oth}}\bigl(\tilde\nu_i + s\,x_{q,k}\,\Delta\tilde\nu_{D,i}\bigr)\Bigr],
  \qquad
  \tau_{\mathrm{oth}}(\tilde\nu) = \sum_{j\neq i,\ \tau_{0,j}\ge10^{-2}} \tau_{0,j}\,\frac{H(a_j,x_j)}{H(a_j,0)},
  \label{eq:overlap}
\end{equation}
with $x_j=(\tilde\nu-\tilde\nu_j)/\Delta\tilde\nu_{D,j}$ and
$\Delta\tilde\nu_{D,j}=\tilde\nu_j b/c$ [cm$^{-1}$]. $\mathcal{O}_i$ approximates
$\int\phi_ie^{-\tau_i}e^{-\tau_{\mathrm{oth}}}\dd\nu\big/\int\phi_ie^{-\tau_i}\dd\nu$, the mean of
$e^{-\tau_{\mathrm{oth}}}$ weighted by line $i$'s own absorption, by a midpoint rule in
the cumulative absorption: each side of the line is cut into three parts that
carry equal shares of its absorption, and each part is represented by the
offset $x_{q,k}$ that halves its share ($1/6$, $1/2$ and $5/6$ of the side). A thin
line thus samples its core and a damped line its wings, hundreds of Doppler
widths out. Three points per side were chosen by test: against an exact
reference the factor on the Solomon rate is 0.86--1.03 with one point per side,
0.99--1.02 with two and 0.98--1.04 with three, and the line-by-line rms falls from
0.03--0.45 to 0.03--0.21 (\dlog{H-006}). Numerically:
\begin{itemize}
\item only lines with $\tau_0\ge10^{-2}$ shade; a shading line's reach is
  $\Delta\tilde\nu_D\max\bigl(12,\sqrt{\tau_0 a/(\sqrt\pi H(a,0)\,\varepsilon)}\bigr)$
  with $\varepsilon=10^{-4}$, the offset where its wing optical depth falls
  to $10^{-4}$;
\item the full Voigt function is used within 12 Doppler widths and the
  Lorentzian wing $\tau_0 a/(\sqrt\pi H(a,0)x^2)$ beyond;
\item lines whose reach exceeds 100\,cm$^{-1}$ (``long'') shade every line
  through their Lorentzian wings, summed as $\sum_j C_j/(\tilde\nu-\tilde\nu_j)^2$
  with four interleaved partial sums; other lines (``short'') are added where
  they reach, within 100\,cm$^{-1}$; lines within 40 Doppler widths of line $i$
  (``near'') are evaluated term by term;
\item the lines are processed in wavenumber order so that the windows only
  move forward; a thin line ($\tau_0<1$) takes the far contribution once at
  its centre; lines that no shading line reaches keep $\mathcal{O}=1$.
\end{itemize}
In \rtn{solveCell} the factor is computed once per cell and path, at the columns
in front of the cell plus half the cell's own (Section~\ref{sec:solvecell}), and
multiplies the cell-averaged $\langle S\rangle$ of Equation~(\ref{eq:cellavg}): a product of
means stands for the mean of the product, adequate where the other lines'
shading changes little across the cell. The result was checked against the
same integral on a 0.05\,cm$^{-1}$ grid across 84\,000--119\,500\,cm$^{-1}$ (every line
with $\tau_0>10^{-3}$ a Voigt profile). The exact factor on the Solomon rate of the
\dlog{H-005} slab's ground-state levels is 0.96, 0.93, 0.86, 0.65, 0.49 and 0.32 at
$N(\mathrm{H_2})=10^{17}$, $10^{18}$, $10^{19}$, $10^{20}$, $3\times10^{20}$ and
$10^{21}$\,cm$^{-2}$; the model reproduces it to 0.2--1.6\% up to
$3\times10^{20}$\,cm$^{-2}$ and to 4.5\% at $10^{21}$\,cm$^{-2}$, and the exact factor is
within 5\% of the overlap term of Equation~(\ref{eq:db96}) (\dlog{H-006}).
Overlap lowered the total pumping through that slab by 33\%.

Not included: a photon absorbed in another line's wing is counted as lost
rather than as a pump of that line (cross-pumping), which underestimates the
pumping of the weaker lines where the band is saturated; and the
H\,\textsc{i} Lyman-series lines, which the MC field does not contain either
(Section~\ref{sec:eta}), do not shade the band. Their damping wings remove
continuum around Ly$\beta$, Ly$\gamma$ and Ly$\delta$ (1026, 973 and 950\,\AA) over
equivalent widths of about 4, 1.4 and 0.7\,\AA{} at
$N(\mathrm{H\,I})=10^{21}$\,cm$^{-2}$, growing as $N(\mathrm{H\,I})^{1/2}$ (an
estimate of this document from the damping-wing curve of growth, i.e.\ a few
per cent of the 912--1110\,\AA{} band); pumping and dissociation behind such
columns are overestimated (Section~\ref{sec:limitations}). The Meudon code
includes them in its exact transfer \citep{LePetit2006}.

\subsection{Escape probabilities}
\label{sec:escape}

The probability that a photon emitted in line $t$ ($u\to l$) at the cell
centre leaves the cloud is (\rtn{h2LineEscape},
\src{h2\_mod.f90:\allowbreak{}1344--1361}; \src{pdr\_mod.f90:\allowbreak{}5821--5835})
\begin{equation}
  \beta_t = \frac{1}{12}\sum_{k=1}^{12} S\bigl(\sigma_t N_{l,k},\,a_t\bigr),
  \label{eq:beta}
\end{equation}
with $N_{l,k}$ the column of the line's lower level from the cell centre to
the edge of the model along line ray $k$ (Section~\ref{sec:lineRays}), the
cell's own half included, from the populations of the previous sweep. With
complete redistribution the photon is emitted with the absorption profile,
so the probability that it crosses a column in one direction,
$\int\phi\,e^{-\tau}\dd x/\int\phi\,\dd x$, is the function $S$ of
Equation~(\ref{eq:SW}); the mean over the twelve HEALPix directions is the mean
over the sphere. The medium is static (no velocity field, so no Sobolev-type
escape), the escape probability is that of the cell centre rather than a mean
over the cell, and each line is treated alone (no overlap in the trapping).
The Meudon code makes the same on-the-spot approximation for the ultraviolet
photons of the \Htwo{} cascades \citep{LePetit2006}.

\subsection{Formation on grains: the nascent distribution}
\label{sec:formdist}

Molecules formed on grains are placed in the $(v,J)$ distribution of
\citet{TakahashiUehara2001}, as Cloudy applies it \citep{Shaw2005}
(\rtn{h2FormDist}, \src{h2\_mod.f90:\allowbreak{}589--639}). With $D$ the dissociation
energy of $X$ from its $(0,0)$ level (36\,118.11\,cm$^{-1}$, 4.478\,eV; their
$D_e=4.48$\,eV), $E_{\mathrm{off}}=0.273$\,eV the energy of $(0,0)$ above the bottom of the
potential well (their zero-point energy $E_{v0}=0.27$\,eV; 0.273\,eV is Cloudy's
value), $E(v,J)$ the level energies above $(0,0)$, and per surface the
parameters of Table~\ref{tab:tu01}:
\begin{gather}
  E_{vm} = X_{\mathrm{vib}}D + E_{\mathrm{off}},\qquad
  E_v = E(v,0)+E_{\mathrm{off}},\\
  E_{\mathrm{dust}}(v) = D\,X_{\mathrm{dust}}\Bigl[1-\frac{E_v-E_{vm}}{D+E_{\mathrm{off}}-E_{vm}}\,\frac{1-X_{\mathrm{dust}}}{2}\Bigr],\qquad
  E_{\mathrm{H_2}}(v) = D+E_{\mathrm{off}}-E_{\mathrm{dust}}(v),\\
  F(v) = e^{-x_v^2},\quad
  x_v = \begin{cases}\dfrac{E_v-E_{vm}}{g_1E_{vm}} & E_v\le E_{vm}\\[6pt]
                     \dfrac{E_v-E_{vm}}{g_2\,(E_{\mathrm{H_2}}(v)-E_{vm})} & E_v>E_{vm}\end{cases},\\
  E_{\mathrm{rot}}(v) = \bigl(E_{\mathrm{H_2}}(v)-E_v\bigr)\frac{X_{\mathrm{rot}}}{X_{\mathrm{rot}}+X_{\mathrm{trans}}},\qquad
  \phi(v,J)\propto g_{\mathrm{ns}}\,\frac{F(v)}{\sum_vF}\,(2J+1)\,\tfrac12\Bigl[e^{-y^2}+e^{-\Delta E/E_{\mathrm{rot}}}\Bigr],
\end{gather}
with $\Delta E=E(v,J)-E(v,0)$, $y=(\Delta E-E_{\mathrm{rot}})/(E_{\mathrm{rot}}/2)$,
$g_{\mathrm{ns}}=3$ (ortho) or 1 (para), $\phi=0$ where $E_{\mathrm{rot}}\le0$ (high $v$;
J.~Takahashi's prescription as quoted in Cloudy's source), and $\phi$
normalised to one. These are equations~1--7 of \citet{TakahashiUehara2001}:
the energy left on the grain $E_{\mathrm{dust}}(v)$ (their eq.~1), the energy that
stays in the molecule (eq.~2) and its share in rotation (eq.~3), and a
vibrational distribution peaked at $E_{vm}$ with different widths below and
above it (eqs.~5--7). The rotational distribution is the mean of their two
models, a Gaussian about $E_{\mathrm{rot}}$ of width $E_{\mathrm{rot}}/2$ (model A,
eq.~12) and a thermal distribution at $E_{\mathrm{rot}}$ (model B, eq.~13); the mean
is applied before one overall normalisation, as in Cloudy. Carbonaceous,
silicate and icy surfaces put the new molecules mainly near $v=2$, 6 and 7
respectively \citep{Shaw2005}. The fraction formed above 4100\,cm$^{-1}$ agrees
with Cloudy's network to four digits (0.9771, \dlog{H-004}).

The energy budget of a formation follows by conservation: $E_{\mathrm{grain}}=\langle
E_{\mathrm{dust}}\rangle_\phi$ is left on the grain,
$E_{\mathrm{kin}}=D-E_{\mathrm{grain}}-\langle E(v,J)\rangle_\phi$ is the translational energy
of the departing molecule, which heats the gas directly, and the internal
energy $\langle E(v,J)\rangle_\phi$ is either radiated in \Htwo{} lines or returned to
the gas by collisional de-excitation, i.e.\ through
$\mathcal{H}_{\mathrm{coll}}$ (Section~\ref{sec:solvecell}). $E_{\mathrm{kin}}$ is 0.834\,eV for
silicate and 1.142\,eV for carbonaceous surfaces, against 3D-PDR's 1.5\,eV
(\dlog{H-013}). $E_{\mathrm{grain}}$ is a diagnostic only (\src{h2Cells.out}); it is not
added to MOCASSIN's grain energy balance.

\begin{table}[H]
\centering
\caption{\citet[their Table~2]{TakahashiUehara2001} parameters as used
(\src{h2\_mod.f90:\allowbreak{}132--139}); the ice set is not used by \rtn{solveCell}.}\label{tab:tu01}
\begin{tabular}{@{}lcccccc@{}}
\toprule
Surface & $X_{\mathrm{vib}}$ & $X_{\mathrm{dust}}$ & $X_{\mathrm{rot}}$ & $X_{\mathrm{trans}}$ & $g_1$ & $g_2$ \\
\midrule
ice & 0.70 & 0.04 & 0.14 & 0.12 & 0.3 & 0.6 \\
silicate & 0.60 & 0.10 & 0.15 & 0.15 & 0.4 & 0.6 \\
carbonaceous & 0.20 & 0.40 & 0.15 & 0.25 & 0.9 & 0.4 \\
\bottomrule
\end{tabular}
\end{table}

In each cell the silicate and carbonaceous distributions are mixed in the
proportion of the chemistry's own formation rate on each material:
$w_{\mathrm{sil}}$ and $w_{\mathrm{gra}}$ are \rtn{H2\_FORMATION\_RATE}
(Section~\ref{sec:formation}) evaluated with only that material's grain area
(\src{pdr\_mod.f90:\allowbreak{}5839--5876}). A cell with no grain formation (no dust,
or $T_d=0$) takes a Boltzmann distribution at $T_e$, i.e.\ no formation pumping,
and is counted. Formation through H$^-$ and \Htwo$^+$, which populates specific
levels, is not distributed over the levels: every molecule the closure forms
enters in $\phi$.

\subsection{The per-cell algorithm}
\label{sec:solvecell}

\rtn{solveCell} (\src{pdr\_mod.f90:\allowbreak{}5665--5954}) does, for one cell:
\begin{enumerate}
\item \textbf{Field} (:5667--5695): $b$ from $T_e$; $J_\nu^{\mathrm{dir}}$,
  $J_\nu^{\mathrm{dif}}$ per bin and $\eta^{\mathrm{dir}}$,
  $\eta^{\mathrm{dif}}$ per line (Equation~\ref{eq:eta}).
\item \textbf{Cell geometry} (:5697--5706): cell widths, $\ell_{\mathrm{iso}}$,
  and the cell's own column for the diffuse paths from its populations of the
  last sweep, $\Delta N^{\mathrm{iso}}_l = n_l n(\mathrm{H_2})\ell_{\mathrm{iso}}$.
\item \textbf{Diffuse shielding} (:5712--5802), if any line has
  $\eta^{\mathrm{dif}}>0$: \rtn{traceDiffuseSites} gives the sites and their
  weights $w$ (Section~\ref{sec:sites}); walking the stored rays again, each
  site gets a column per level,
  \begin{equation}
    N_{l,\mathrm{site}} = N^\star_l(P) + n_l(P)\,n(\mathrm{H_2},P)\,\ell_{\mathrm{own}}(P)
        + \sum_{\mathrm{between}} n_l\,n(\mathrm{H_2})\,\ell + \tfrac12\,n_l(P)\,n(\mathrm{H_2},P)\,\ell_P
  \end{equation}
  (the level-resolved form of Equation~\ref{eq:twoleg}, with the source-ward
  level columns $N^\star_l$ of \rtn{starColumns}). Evaluating $S$ for every
  line at every site would cost $32\,453\times N_{\mathrm{site}}$ look-ups,
  and a cell has hundreds of sites, so the sites are binned:
  \begin{itemize}
  \item four groups by total column $\sum_l N_{l,\mathrm{site}}$, with edges
    $10^{17}$, $10^{19}$ and $3\times10^{20}$\,cm$^{-2}$;
  \item within a group and for each level, a histogram of the weights in
    $\log_{10}N_{l,\mathrm{site}}$, 0.1\,dex bins from $10^9$\,cm$^{-2}$
    (150 bins, the end bins catching the columns beyond), plus a ``thin'' weight for sites where
    $\sigma_{\max,l}N_l<10^{-4}$, $\sigma_{\max,l}$ being the largest
    cross section of the lines out of $l$;
  \item the overlap factor (Equation~\ref{eq:overlap}) once per group with
    more than $10^{-4}$ of the total weight, at the group's weight-mean
    level columns plus $\frac12\Delta N^{\mathrm{iso}}$ (none for the
    $<10^{17}$\,cm$^{-2}$ group);
  \item per line, the diffuse shielding factor
    \begin{equation}
      \langle S\rangle^{\mathrm{dif}}_i = \frac{1}{\sum w}\sum_{g}\mathcal{O}_{i,g}
      \Bigl[w^{\mathrm{thin}}_{l,g}\,\langle S\rangle(0) + \sum_{\mathrm{bins}} w_{l,g,\mathrm{bin}}\,
      \langle S\rangle\bigl(\sigma_i N_{\mathrm{bin}}\bigr)\Bigr],
    \end{equation}
    each $\langle S\rangle$ averaged through the cell's own column
    $\sigma_i\Delta N^{\mathrm{iso}}_l$ by Equation~(\ref{eq:cellavg}).
  \end{itemize}
  The binning replaces each site's column by its bin centre (an error of at
  most $\pm0.05$\,dex in $N_l$) and each group's overlap by that at the group's
  mean columns. If there is diffuse light but no site (no dust on any of the
  twelve rays) the diffuse part rides the direct path (:5799--5802).
\item \textbf{Direct path} (:5804--5819): source-ward level columns
  $N^\star_l$ and the chord $\ell_{\mathrm{dir}}$.
\item \textbf{Escape} (:5821--5835): $\beta_t$ over the twelve line rays
  (Equation~\ref{eq:beta}).
\item \textbf{Formation} (:5839--5876): mixture weights, $\phi$, $E_{\mathrm{kin}}$,
  $E_{\mathrm{grain}}$.
\item \textbf{Solve} (:5879--5914). \rtn{h2LineShieldPrep} stores, for the
  direct column, each line's $\tau_0$, $W(\tau_0)$ and damping index. Then up
  to 12 passes: colliders from the current ortho fraction; the cell's own
  direct column $\Delta N^{\mathrm{dir}}_l = n_l n(\mathrm{H_2})\ell_{\mathrm{dir}}$
  from the current populations; on the first pass only, the direct overlap factor at
  $N^\star+\frac12\Delta N^{\mathrm{dir}}$; the direct shielding
  $\langle S\rangle^{\mathrm{dir}}_i\mathcal{O}^{\mathrm{dir}}_i$
  (\rtn{h2LineShieldEval}); the shielded occupation numbers
  \begin{equation}
    \eta_i = \eta^{\mathrm{dir}}_i\,\langle S\rangle^{\mathrm{dir}}_i\mathcal{O}^{\mathrm{dir}}_i
           + \eta^{\mathrm{dif}}_i\,\langle S\rangle^{\mathrm{dif}}_i
  \end{equation}
  (or $(\eta^{\mathrm{dir}}_i+\eta^{\mathrm{dif}}_i)$ times the direct factor
  when there is no site); and the level balance (\rtn{h2Solve}). The loop stops
  when the ortho fraction changes by less than $10^{-4}$; a cell that had
  settled in the previous sweep needs one pass, and the mean was 1.0--1.6 passes
  per cell (\dlog{H-011}).
\item \textbf{Products} (:5917--5952): $k_{\mathrm{Sol}}=\sum n_lk_l$;
  the absorption rate $\sum n_l\sum_uP_{lu}$; for comparison the
  chemistry's own single-line rate on the same field and columns
  (\rtn{h2SolChem}, 3D-PDR's \rtn{H2PDRATE}); diagnostics; and three
  heating terms per molecule: the atoms' kinetic energy from dissociation
  $\sum_l n_l H^{\mathrm{dis}}_l$, the net heating by collisions within $X$
  \begin{equation}
    \mathcal{H}_{\mathrm{coll}} = hc\sum_c\sum_{u>l} n_c k^c_{ul}
       \Bigl(n_u - n_l\frac{g_u}{g_l}e^{-hc(E_u-E_l)/kT}\Bigr)(E_u-E_l)
       \quad[\mathrm{erg\,s^{-1}}]
    \label{eq:hcoll}
  \end{equation}
  (\rtn{h2CollHeat}, \src{h2\_mod.f90:\allowbreak{}1425--1448}; Cloudy's
  \texttt{HeatDexc}), and the kinetic energy per molecule formed,
  $E_{\mathrm{kin}}$. $\mathcal{H}_{\mathrm{coll}}$ is the energy that collisions take out
  of the levels populated above their Boltzmann ratio at $T$ by pumping and
  formation; it is negative, i.e.\ \Htwo{} line cooling, where the gas is
  hotter than the level populations. A pump deposits in the gas only what
  $\mathcal{H}_{\mathrm{coll}}$ recovers from the levels it feeds, and a dissociation
  the kinetic energy of the atoms, so the three terms do not count any
  energy twice.
\end{enumerate}

\subsection{Sweeps, parallel strategy and convergence}
\label{sec:sweeps}

\rtn{h2LevelModel} reads the data once and, on every call:
\begin{itemize}
\item takes $n(\mathrm{H_2})=n_{\mathrm H}x(\mathrm{H_2})$ of every cell
  from the chemistry; on the first call it starts every cell from Boltzmann
  populations at $T_e$, or, on a warm start, from \src{output/\allowbreak{}h2Pop.bin}
  if that file describes the same grid and level set (\rtn{h2PopResume},
  \src{pdr\_mod.f90:\allowbreak{}5961--6012}, \dlog{H-020}); later calls start from the last
  call's populations. A cold start on a converged model overshoots the
  dissociations of its first sweep by 38\% (\dlog{H-016}); a resumed one is
  within 0.6\% of the checkpointed run (\dlog{H-020});
\item for each sweep, rebuilds the source-ward level columns of every cell,
  $N^\star_l=\sum n_l\,n(\mathrm{H_2})\,\ell$ along \rtn{rayCells} to each
  source, weighted over sources by $r_s^{-2}$ and excluding the cell itself
  (\rtn{starColumns}, \src{:\allowbreak{}5621--5659}); the other cells' populations in all
  columns are those of the previous sweep, so a sweep is a Jacobi iteration
  over the cells, whose result does not depend on the order in which the
  cells are solved;
\item distributes the cells dynamically: the ranks take cells one at a time
  from a shared counter in an MPI-3 one-sided window allocated on rank~0
  (\verb|MPI_Win_allocate|, \verb|MPI_Fetch_and_op|, \src{:\allowbreak{}5398--5426}); the
  first sweep takes them in a fixed pseudo-random order (a Fisher--Yates shuffle
  driven by a linear congruential generator), later sweeps in decreasing cost
  of the previous sweep (\rtn{orderCells}, \src{:\allowbreak{}5538--5559}), so that cores of
  different speed finish together; a cell's result does not depend on the rank
  that solves it;
\item reduces populations, rates and heating terms with \verb|MPI_Allreduce|
  (in single precision);
\item stops when the luminosities of 1--0\,S(1), 2--1\,S(1), 0--0\,S(0) and
  0--0\,S(1), summed over the grid in the optically thin limit, all change by
  less than 1\% between sweeps (\src{:\allowbreak{}5459--5495}), at most \rtn{nSweepMax}
  sweeps.
\end{itemize}
Inside the MC iteration \rtn{nSweepMax}$=1$: one sweep per iteration, the
level model converging together with the rest of the model. On the finished
model \rtn{nSweepMax}$=12$ (Section~\ref{sec:convergence}). On \texttt{pn150k}
($41^3$, 35\,384 cells, 14 ranks) a call took 3736\,s before and 1474\,s after the
optimisations of \dlog{H-011}; the overlap calculation was about half of the
timed work (\dlog{H-007}).
